\paragraph{Letting renters decline a higher ask.}
After observing the 860 rent increases, we tested an optional rule. A renter
with a provisional unit can decline to apply for any displayed alternative
with a higher asking rent. Their own unit stays on their effective list.
Deleting other edges preserves the stability of the original assignment.
Apply Corollary~\ref{cor:guard-rank} to this restricted original market
and retain the same guarded higher-income lists. A consenting renter's
new unit is weakly preferred and costs no more than their original unit.
The rule changes which units the renter consents to pursue. The platform
still displays the full quota to everyone.

We simulated take-up by all 15,847 original lower-income recipients,
using the original catalogue asking rents.
All retained their rank floor, none faced a rent increase, and 1,345 strictly
improved. Compared with protection without the rent veto, 967 people gave
up a strict improvement and returned to their original unit, and 72 others gained
an improvement for the first time. Lower-income applicants without a provisional
offer also experienced the effects. In this group 4,181 gained a placement under the veto,
compared with 4,422 without it. Across the two seasons and designs the higher
group gained 64 placements and total placements fell by 177. In the enriched
markets the veto reduced the lower group's match rate by 0.61 percentage
points in summer (paired 95\% interval $[-0.75,-0.46]$) and 0.65 in winter
($[-0.80,-0.49]$). Actual take-up and willingness to pay for a preferred
unit remain open empirical questions.
